% Options for packages loaded elsewhere
\PassOptionsToPackage{unicode}{hyperref}
\PassOptionsToPackage{hyphens}{url}
%
\documentclass[
  11pt,
]{article}
\usepackage{amsmath,amssymb}
\usepackage{setspace}
\usepackage{iftex}
\ifPDFTeX
  \usepackage[T1]{fontenc}
  \usepackage[utf8]{inputenc}
  \usepackage{textcomp} % provide euro and other symbols
\else % if luatex or xetex
  \usepackage{unicode-math} % this also loads fontspec
  \defaultfontfeatures{Scale=MatchLowercase}
  \defaultfontfeatures[\rmfamily]{Ligatures=TeX,Scale=1}
\fi
\usepackage{lmodern}
\ifPDFTeX\else
  % xetex/luatex font selection
\fi
% Use upquote if available, for straight quotes in verbatim environments
\IfFileExists{upquote.sty}{\usepackage{upquote}}{}
\IfFileExists{microtype.sty}{% use microtype if available
  \usepackage[]{microtype}
  \UseMicrotypeSet[protrusion]{basicmath} % disable protrusion for tt fonts
}{}
\makeatletter
\@ifundefined{KOMAClassName}{% if non-KOMA class
  \IfFileExists{parskip.sty}{%
    \usepackage{parskip}
  }{% else
    \setlength{\parindent}{0pt}
    \setlength{\parskip}{6pt plus 2pt minus 1pt}}
}{% if KOMA class
  \KOMAoptions{parskip=half}}
\makeatother
\usepackage{xcolor}
\usepackage[margin=0.85in]{geometry}
\setlength{\emergencystretch}{3em} % prevent overfull lines
\providecommand{\tightlist}{%
  \setlength{\itemsep}{0pt}\setlength{\parskip}{0pt}}
\setcounter{secnumdepth}{-\maxdimen} % remove section numbering
\usepackage{amsmath,amssymb,microtype}
\usepackage{xurl}
\setlength{\emergencystretch}{2em}
\widowpenalty=10000
\clubpenalty=10000
\displaywidowpenalty=10000
\AtBeginDocument{\hypersetup{pdfsubject={Explicit fixed-dimensional simplex rigidity with a stronger inherited planar comparison}}}
\pdfinfoomitdate=1
\pdftrailerid{}
\ifLuaTeX
  \usepackage{selnolig}  % disable illegal ligatures
\fi
\usepackage{bookmark}
\IfFileExists{xurl.sty}{\usepackage{xurl}}{} % add URL line breaks if available
\urlstyle{same}
\hypersetup{
  pdftitle={Effective simplex rigidity for the projection cone invariant},
  pdfauthor={mxym},
  hidelinks,
  pdfcreator={LaTeX via pandoc}}

\title{Effective simplex rigidity for the projection cone invariant}
\author{mxym}
\date{Supplement to entry 005, 7 October 2026}

\begin{document}
\maketitle

\setstretch{1.04}
\subsection{Abstract}\label{abstract}

For every fixed dimension \(d\ge2\), we give a fully explicit power
modulus for the lower equality case of the projection-cone invariant. If
\(e(K)=a(K)-1/(d+1)\) and \(S\) is any maximum-volume simplex inscribed
in \(K\), with centroid \(z\), then
\[S\subseteq K\subseteq z+(1+G_d e(K)^{\theta_d})(S-z),
\qquad \theta_d=\frac1{d(d^2+2d+2)}.\] All constants are specified. The
proof works for arbitrary convex bodies and cone-volume laws, including
nonatomic laws. Its constants are extremely poor and no optimality or
novelty is asserted. In dimension two the exact identity
\(a(K)=2|K|/|K-K|\) connects the question to classical Rogers--Shephard
stability: an estimate in Böröczky's work already implies the much
stronger linear bound \(d_{BM}(K,T_2)\le1+576e(K)\). Only the simplex
endpoint is treated here; version 4 concerns the symmetric product
endpoint.

\subsection{1. Statement and scope}\label{statement-and-scope}

A convex body is a compact convex set with nonempty interior. For
\(K\subset\mathbb R^d\), let \(|K|\) denote volume and let \(\Pi K\) be
its projection body, normalized by
\[h_{\Pi K}(u)=|K\mid u^\perp|_{d-1}\qquad (|u|=1).\] Put
\[R(K)=\frac{|\Pi K|}{|K|^{d-1}},\qquad
a(K)=\left(\frac d{d+1}\right)^d
\frac{R(\mathcal P K)}{R(K)}-1,\] where \(\mathcal P K\) is a pyramid
over \(K\). Versions 2--3 establish that \(a\) is a continuous affine
invariant, give its cone-volume representation recalled below, and prove
\(a(K)\ge1/(d+1)\) with equality exactly for simplices. We use those
identities and standard Cauchy and mixed-volume formulas. The full
quantitative argument is supplied here.

For any maximum-volume inscribed simplex \(S\) with centroid \(z\),
define \[E(K,S)=\inf\{t\ge0:K\subseteq z+(1+t)(S-z)\}.\] The
multiplicative affine Banach--Mazur distance from \(K\) to a simplex is
the infimum of dilation factors for a sandwich between an inscribed
simplex and a translate of its dilate. In particular it is at most
\(1+E(K,S)\).

\textbf{Explicit constants.} For \(d\ge2\), define, in this order,
\[\begin{aligned}
R_0&=d(d+1),& b&=\frac1{4(dR_0)^d},&M&=\frac2b,\\
L&=d2^d,&Z&=M(2d+3),&C&=4d^2R_0(M+1),\\
p&=d^2+2d+2,&\theta_d&=\frac1{dp},&
\chi&=b^d\left(\frac{b^d}{12dL}\right)^{d(d+1)}.
\end{aligned}\] Next set \[\begin{aligned}
\eta_0&=\min\left\{\frac1{2(d+1)},\frac b{4d},\frac1Z,
                  \frac1{Z(8MdC)^d}\right\},\\
e_0&=\frac{\chi\eta_0^p}{2(d+1)},\\
A_d&=4MdC Z^{1/d}\left(\frac{2(d+1)}\chi\right)^{\theta_d},\\
G_d&=\max\{A_d,(R_0-1)e_0^{-\theta_d}\}.
\end{aligned}\] These are explicit numbers depending only on \(d\);
\(\chi,\eta_0,e_0\) are rational.

\textbf{Theorem 1 (effective simplex rigidity).} For every convex body
\(K\subset\mathbb R^d\), \(d\ge2\), and every maximum-volume inscribed
simplex \(S\),
\[E(K,S)\le G_d e(K)^{\theta_d},\qquad e(K)=a(K)-\frac1{d+1}.\] If
\(0\le e(K)\le e_0\), the sharper stated constant applies:
\[E(K,S)\le A_d e(K)^{\theta_d},\qquad A_de_0^{\theta_d}\le\frac12.\]
Consequently
\[\delta(d,\varepsilon)=\min\{e_0,(\varepsilon/A_d)^{dp}\}\] is an
explicit choice in version 3's qualitative containment theorem. In
dimension one \(K\) is itself an interval and \(E(K,S)=0\).

The first exponents are \(1/20,1/51,1/104\) in dimensions \(2,3,4\).
Already for \(d=2\) the displayed local threshold satisfies
\(10^{-279}\le e_0<10^{-278}\). The theorem supplies an effective
modulus, not a practical numerical estimate. Section 8 explains the
stronger inherited planar result. No dimension-independent conclusion
follows.

\subsection{2. Normalization and the exact cancellation
defect}\label{normalization-and-the-exact-cancellation-defect}

Fix an arbitrary maximum simplex \(S\). Map it affinely to a regular
simplex \(\Delta\) with centroid zero and inradius one. Its vertices
\(v_0,\ldots,v_d\) have norm \(d\). If \(\alpha_i(x)\) are barycentric
coordinates relative to \(\Delta\), replacing \(v_i\) by \(x\in K\)
multiplies volume by \(|\alpha_i(x)|\). Maximality gives
\(|\alpha_i(x)|\le1\), hence
\[B_2^d\subseteq\Delta\subseteq K\subseteq R_0 B_2^d.\tag{2.1}\] The
deficit is unchanged by this normalization.

Write \(h=h_K\) and let \(S_K\) denote surface area measure. Push the
probability measure \[\frac{h(u)}{d|K|}\,dS_K(u)\] forward by
\(u\mapsto u/h(u)\), obtaining \(\nu\). This centered law is supported
on \(\partial K^\circ\subseteq B_2^d\). For independent samples \(X_i\)
of law \(\nu\), put \[A=\mathbb E|\det(X_1,\ldots,X_d)|,\qquad
B=\mathbb E|\det((X_1,1),\ldots,(X_{d+1},1))|.\] The cone-volume
representation from version 3 reads \[a(K)=\frac B{(d+1)A},\qquad
D:=B-A=(d+1)Ae(K)\le(d+1)e(K),\tag{2.2}\] because \(0<A\le1\). For every
unit vector \(u\), centering and Cauchy's projection formula give
\[\int(-\langle u,x\rangle)_+\,d\nu(x)
=\frac{|K\mid u^\perp|_{d-1}}{d|K|}
\ge\frac{(2/d)^{d-1}}{d(2R_0)^d}=2b.\tag{2.3}\] Here the projection
contains the unit ball in dimension \(d-1\), which contains a cube of
side \(2/d\).

For an ordered base tuple \(\mathbf x=(x_1,\ldots,x_d)\) define
\[F_{\mathbf x}(y)=\det((x_1,1),\ldots,(x_d,1),(y,1))
=\det(x_1-y,\ldots,x_d-y),\] and
\(P(\mathbf x)=\int(F_{\mathbf x})_+\,d\nu\),
\(N(\mathbf x)=\int(-F_{\mathbf x})_+\,d\nu\). Centering gives
\[P(\mathbf x)-N(\mathbf x)=\det(x_1,\ldots,x_d).\] Thus the exact
cancellation identity is
\[D=2\int\min\{P(\mathbf x),N(\mathbf x)\}\,d\nu^{\otimes d}(\mathbf x).\tag{2.4}\]
This identity includes horizontally singular base tuples.

If two \((d+1)\)-tuples in \(B_2^d\) have corresponding points at
distance at most \(r\), their affine determinants differ by at most
\[Lr=d2^dr.\tag{2.5}\] Indeed, write it using the \(d\) difference
columns above and telescope one column at a time. Each old or new column
has norm at most two and each change has norm at most \(2r\).

\subsection{3. Concentration on heavy
cells}\label{concentration-on-heavy-cells}

Fix \(0<\eta\le\eta_0\) and set \[r=\frac{\eta b^d}{4L},\qquad
\tau=\eta\left(\frac r{3d}\right)^d,\qquad
\gamma=\eta b^d.\tag{3.1}\] Partition \([-1,1]^d\) into measurable cubes
of diameter at most \(r\), assigning boundaries to single cells.
Dividing each axis into \(\lceil2\sqrt d/r\rceil\) intervals uses at
most \((3d/r)^d\) cells, since \(r\le1\). Call a cell heavy if its
\(\nu\)-mass is at least \(\tau\). The light cells together have mass at
most \(\eta\).

Choose one support point in each heavy cell and let \(X\) be this finite
set. The chosen points need not have positive individual mass. The
function \((-\langle u,x\rangle)_+\) is one-Lipschitz and bounded by one
on the support. Discarding light cells and replacing samples by
representatives therefore gives
\[\sum_{Q\text{ heavy}}\nu(Q)(-\langle u,x_Q\rangle)_+
\ge2b-\eta-r\ge b.\] The weights sum to at most one, so for every unit
\(u\) some representative has \(\langle u,x_Q\rangle\le-b\). Hence
\[bB_2^d\subseteq\operatorname{conv}X.\tag{3.2}\]

Choose a maximum-volume simplex
\(T=\operatorname{conv}(w_0,\ldots,w_d)\) with vertices in \(X\), and
put \[V=|\det((w_0,1),\ldots,(w_d,1))|.\] Then \(V\ge b^d\): the simplex
with vertices \(0,be_1,\ldots,be_d\) lies in \(\operatorname{conv}X\),
and separate convexity of the absolute affine determinant lets each of
its vertices be replaced successively by a point of \(X\) without
decreasing the determinant. At this stage we do not assume that
\(0\in T\).

Write \(x=\sum_i\alpha_iw_i\) for \(x\in X\), with \(\sum_i\alpha_i=1\).
There are two possible robust witnesses:

\begin{enumerate}
\def\labelenumi{\arabic{enumi}.}
\tightlist
\item
  If some \(\alpha_i\le-\eta\), use the \(d\) vertices other than
  \(w_i\) as base. The determinant values at \(w_i\) and \(x\) have
  opposite signs and magnitudes at least \(\eta V\).
\item
  If \(\alpha_i,\alpha_j\ge\eta\) for two distinct indices, use \(x\)
  and the \(d-1\) original vertices other than \(w_i,w_j\) as base.
  Expansion of \((x,1)\) gives opposite determinant values at
  \(w_i,w_j\), of magnitudes \(\alpha_jV\) and \(\alpha_iV\).
\end{enumerate}

In each case the \(d\) base representatives and the two tested
representatives lie in \(d+2\) distinct heavy cells. Since
\(Lr=\gamma/4\) and \(\eta V\ge\gamma\), inequality (2.5) preserves both
signs with margin at least \(\gamma/2\) for every tuple of actual
support points in the corresponding cells. For each sampled base tuple,
both integrals in (2.4) are at least \(\tau\gamma/2\). The ordered
product of base cells has mass at least \(\tau^d\). Consequently either
witness forces \[D\ge\tau^{d+1}\gamma=\chi\eta^p.\tag{3.3}\] This
integrates over full positive-measure cells, not over the selected
points.

Assume now that \[D<\chi\eta^p.\tag{3.4}\] Both witnesses are
impossible. Every representative has all coefficients greater than
\(-\eta\) and at most one coefficient at least \(\eta\). There is
exactly one such coefficient since \(\eta\le1/[2(d+1)]\). The other
\(d\) coefficients have absolute value at most \(\eta\), whence
\[\operatorname{dist}(x,\{w_0,\ldots,w_d\})\le2d\eta
\qquad(x\in X).\] Combining this with (3.2), support functions and
\(2d\eta\le b/2\) proves
\[\frac b2 B_2^d\subseteq T\subseteq B_2^d.\tag{3.5}\] Thus zero is now
known to be interior to \(T\). Assign a sample in each heavy cell to its
nearby vertex and assign light-cell samples arbitrarily. The resulting
expected distance satisfies
\[\int\operatorname{dist}(x,\{w_0,\ldots,w_d\})\,d\nu(x)
\le r+2d\eta+2\eta\le(2d+3)\eta.\tag{3.6}\]

\subsection{4. An enclosing simplex and a volume
estimate}\label{an-enclosing-simplex-and-a-volume-estimate}

Let \[P=T^\circ=\{y:\langle w_i,y\rangle\le1\text{ for all }i\}.\] Each
\(w_i\) belongs to \(\partial K^\circ\). By (3.5), \(P\) is a bounded
simplex and \[K\subseteq P\subseteq MB_2^d.\tag{4.1}\] Every \(w_i\) is
a vertex of \(T\), so \(h_P(w_i)=1\). Also \(h_K(x)=1\) on
\(\operatorname{supp}\nu\). The support function \(h_P\) is
\(M\)-Lipschitz, and (3.6) implies
\[0\le\int(h_P(x)-1)\,d\nu(x)\le Z\eta.\tag{4.2}\] The left-hand side
equals \[\frac1{d|K|}\int h_P(u)\,dS_K(u)-1
=\frac{V(K[d-1],P)}{|K|}-1.\] Minkowski's first inequality therefore
gives \[|P|^{1/d}\le|K|^{1/d}(1+Z\eta).\] Since \(Z\eta\le1\),
\[|P|-|K|\le(2R_0)^d d2^{d-1}Z\eta.\tag{4.3}\] The conversion uses the
constructed circumscribed simplex directly; it does not require an
inverse-continuity constant for surface area measures.

\subsection{5. A cap estimate}\label{a-cap-estimate}

Put \(s=d_H(K,P)\). If \(s>0\), choose a farthest point \(q\in P\) from
\(K\) and the unit direction \(u\) from its metric projection onto
\(K\), so \[\langle u,q\rangle-h_K(u)=s.\] Set \(t=s/[2(M+1)]\). We have
\(s\le M\), hence \(0<t<1\). The ball \[ (1-t)q+tB_2^d\] lies in
\(\operatorname{conv}(q,B_2^d)\subseteq P\). Its lowest \(u\)-coordinate
is at least \(h_K(u)+s/2\), because \(\langle u,q\rangle\le M\). It is
therefore disjoint from \(K\). Its inscribed cube has side \(2t/d\),
giving \[|P|-|K|\ge\left(\frac{s}{d(M+1)}\right)^d.\] Together with
(4.3), this yields, also when \(s=0\),
\[s\le2dR_0(M+1)[d2^{d-1}]^{1/d}(Z\eta)^{1/d}
\le C(Z\eta)^{1/d}.\tag{5.1}\] The last term in the definition of
\(\eta_0\) ensures \[s\le\frac1{8Md}.\tag{5.2}\]

\subsection{6. Returning to the original maximum
simplex}\label{returning-to-the-original-maximum-simplex}

The conclusion must concern the original arbitrarily chosen \(\Delta\),
not just the simplex \(P\) constructed from the law. Since \(h_P\ge1\)
on the unit sphere and \(h_K\ge h_P-s\),
\[ (1-s)P\subseteq K\subseteq P.\] Maximality of \(\Delta\) then gives
\[\frac{|\Delta|}{|P|}\ge(1-s)^d\ge1-ds.\] Put \(\delta=ds\le1/8\). The
barycentric coordinates of the \(d+1\) vertices of \(\Delta\) relative
to \(P\) form a nonnegative column-stochastic matrix \(W\) satisfying
\[|\det W|=\frac{|\Delta|}{|P|}\ge1-\delta.\tag{6.1}\] Each column has
Euclidean norm at most one. Hadamard's inequality forces the squared
norm of each column to be at least \((1-\delta)^2\), so its largest
entry is at least \(1-2\delta\).

The dominant rows are distinct. Indeed, the permanent is the probability
that independent categorical draws from the columns choose distinct
rows. A common dominant row in two columns would imply
\[|\det W|\le\operatorname{per}W
\le1-(1-2\delta)^2\le4\delta<1-\delta,\] contrary to (6.1). Relabel the
vertices of \(P\) accordingly. Each vertex of \(\Delta\) is within
\[2\delta\operatorname{diam}P\le4M\delta\] of its matched vertex of
\(P\). Taking convex combinations proves \[d_H(\Delta,P)\le4Md s.\]
Because \(B_2^d\subseteq\Delta\) and \(\Delta\) has centroid zero,
\[K\subseteq P\subseteq\Delta+4Md sB_2^d
\subseteq(1+4Md s)\Delta.\] We have proved the parameterized estimate
\[E(K,\Delta)\le4MdC(Z\eta)^{1/d}\tag{6.2}\] whenever
\(0<\eta\le\eta_0\) and \(D<\chi\eta^p\).

\subsection{7. Completing the explicit
modulus}\label{completing-the-explicit-modulus}

If \(0<e=e(K)\le e_0\), choose
\[\eta=\left(\frac{2(d+1)e}{\chi}\right)^{1/p}.\] Then \(\eta\le\eta_0\)
and (2.2) gives \[D\le(d+1)e=\frac\chi2\eta^p<\chi\eta^p.\] Substituting
in (6.2) proves \(E(K,\Delta)\le A_de^{\theta_d}\). At \(e=0\), the
parameterized bound holds for every sufficiently small positive
\(\eta\); letting \(\eta\downarrow0\) proves the same assertion. No
equality conclusion is obtained by merely passing to a limit of strict
inequalities.

At the threshold, \[A_de_0^{\theta_d}=4MdC(Z\eta_0)^{1/d}\le\frac12.\]
For \(e>e_0\), (2.1) gives \(K\subseteq R_0\Delta\), hence
\[E(K,\Delta)\le R_0-1
\le(R_0-1)e_0^{-\theta_d}e^{\theta_d}.\] This proves the global constant
\(G_d\). Undo the initial affine map: the centroid goes back to \(z\),
and the homotheties become those about \(z\). The original maximum
simplex was arbitrary and all constants depend only on \(d\), so the
result holds for every such simplex. This completes the proof of Theorem
1.

\subsection{8. The planar identity and stronger inherited
stability}\label{the-planar-identity-and-stronger-inherited-stability}

\textbf{Proposition 2 (planar reduction).} Every planar convex body
satisfies \[a(K)R(K)=2,\qquad R(K)=\frac{|K-K|}{|K|}.\]

\textbf{Proof.} First take a polygon with the origin in its interior.
List its edges counterclockwise. Let \(u_i\) be the edge length times
outward unit normal, \(b_i\) the edge length times support number, and
\(w_i=(u_i,b_i)\). Then \(b_i>0\) and \(u_i/b_i\) are the
counterclockwise polar vertices. Hence for \(i<j<k\) all lifted
determinants have positive sign, and the numerator in the polytope
formula for \(a\) is \[\Sigma=\sum_{i<j<k}\det(w_i,w_j,w_k).\] Let
\(T_0=\sum_{i<j}\det(u_i,u_j)\). On expanding \(\Sigma\), the
coefficient of \(b_k\) is \[\sum_{i<j<k}\det(u_i,u_j)
-\sum_{i<k<j}\det(u_i,u_j)
+\sum_{k<i<j}\det(u_i,u_j).\] It equals \(T_0\). To verify this, put
\(U=\sum_{i<k}u_i\) and \(W=\sum_{j>k}u_j\). Closure gives \(W=-U-u_k\),
so \[\det(U,W)=-\det(U,u_k),\qquad
\det(u_k,W)=\det(U,u_k),\] which cancels the difference between the two
expressions. The vectors \(u_i\) are the oriented edges rotated by a
fixed quarter-turn; the signed area formula gives \(T_0=2|K|\). Also
\(\sum_i b_i=2|K|\). Thus \(\Sigma=4|K|^2\).

The denominator in the polytope formula from version 2 is
\(2|K|\,|\Pi K|\), so \(a(K)=2|K|/|\Pi K|\). In the plane the support
function of \(\Pi K\) is the width in a perpendicular direction;
consequently \(\Pi K\) is a quarter-turn of \(K-K\). This proves both
identities for polygons. Polygonal approximation and continuity extend
them to every planar convex body.

Put \(e=a(K)-1/3\) and \(\rho=1-R(K)/6\). Proposition 2 gives
\[\rho=\frac{3e}{1+3e}.\] Theorem 1 of Böröczky {[}B{]} already gives a
linear Banach--Mazur stability estimate. More sharply, Section 8 of the
inspected author-hosted text gives \(R(K)\le6-t/32\) when
\(d_{BM}(K,T_2)=1+t\). Hence the inherited consequence is
\[\boxed{d_{BM}(K,T_2)-1\le192\rho
=\frac{576e}{1+3e}\le576e.}\] This prior planar result is stronger in
exponent than Theorem 1's planar Banach--Mazur consequence. We do not
present planar stability as a new theorem. The reduction to a difference
body is specific to dimension two; no such formula is asserted in higher
dimensions.

\subsection{9. Relation to version 4 and
limitations}\label{relation-to-version-4-and-limitations}

Version 4 classifies the centrally symmetric upper endpoint:
\(a(K)=1/2\) holds exactly, up to an invertible linear map, for
Cartesian products of centrally symmetric bodies of dimensions one and
two. It also gives a fixed-dimensional qualitative stability statement
for that class. The present theorem treats the lower deficit
\(a(K)-1/(d+1)\), whose equality class consists of simplices. It does
not produce a quantitative modulus for \(1/2-a(K)\), and it does not
determine the spectral supremum discussed in versions 2--3.

No boundary smoothness or atom assumption occurs in the proof. The
quantitative witness uses \(d+2\) positive-mass cells; all determinant
signs are made uniform on those cells before integration. The losses in
the exponent come from the covering, the \(d+1\) cell-mass factors, and
the cap estimate. No sharpness or favorable dependence on dimension is
claimed. The previously documented family \(T_m\times T_m\) remains a
barrier to a dimension-independent additive-defect modulus.

Quantitative random-simplex stability {[}AB{]}, cone-volume stability
{[}BH{]}, and inverse stability for surface area measures {[}HS{]} are
established subjects. The particular cone-law deficit and elementary
enclosing-simplex argument are the formulation proved here, not evidence
that it lacks antecedents. Saroglou's cone/projection-body work {[}S{]}
is especially relevant; only its publisher abstract and author listing
were available in this comparison, so formula-by-formula overlap remains
unchecked. This is a bounded literature comparison, not a priority or
best-known-rate claim.

The standard-library exact checks accompany the written proof. They
verify the displayed constants in dimensions 2 through 8, 400
determinant perturbations, 400 barycentric matrices, 8,320 nonatomic
box-corner sign tests, and 200 polygon identities. Finite checks are
regression evidence, not a proof of the general theorem. Independent
model audits are not external human peer review or proof-assistant
formalization.

\clearpage

\subsection{References and provenance}\label{references-and-provenance}

\begin{itemize}
\tightlist
\item
  {[}005v2{]} mxym, \emph{Projection-volume calculus for joins and
  Cartesian products}, entry 005, version 2, with its spectral
  supplement. Definitions and cone/product formulas.
  \href{https://github.com/mxym/math/tree/dd969ea280508a15995fbe83a80cac0f6d227f73/preprints/005-simplex-product-optimum/v2}{Repository
  source}.
\item
  {[}005v3{]} mxym, \emph{Random-determinant rigidity, sharp symmetric
  cone bounds, and spectral nonattainment}, entry 005, version 3.
  Cone-law representation and qualitative lower-end rigidity.
  \href{https://github.com/mxym/math/tree/dd969ea280508a15995fbe83a80cac0f6d227f73/preprints/005-simplex-product-optimum/v3}{Repository
  source}.
\item
  {[}005v4{]} mxym, \emph{Equality in the symmetric projection-cone
  bound}, entry 005, version 4.
  \href{https://github.com/mxym/math/blob/dd969ea280508a15995fbe83a80cac0f6d227f73/preprints/005-simplex-product-optimum/v4/paper.md}{Pinned
  source}.
\item
  {[}B{]} K. Böröczky Jr., \emph{The stability of the Rogers--Shephard
  inequality and of some related inequalities}, Advances in Mathematics
  190 (2005), 47--76, DOI
  \href{https://doi.org/10.1016/j.aim.2003.11.015}{10.1016/j.aim.2003.11.015}.
  \href{https://www.renyi.hu/~carlos/rogerstab.pdf}{Inspected
  author-hosted text}, Theorem 1 and Section 8, printed pages 1 and 20.
  This hosted version carries a later submission timestamp.
\item
  {[}AB{]} G. Ambrus and K. J. Böröczky, \emph{Stability results for the
  volume of random simplices}, American Journal of Mathematics 136
  (2014), 833--857, DOI 10.1353/ajm.2014.0030.
  \href{https://www.renyi.hu/~carlos/randomsimplex.pdf}{Author-hosted
  text}.
\item
  {[}HS{]} D. Hug and R. Schneider, \emph{Hölder continuity for support
  measures of convex bodies}.
  \href{https://arxiv.org/abs/1501.06214}{Primary preprint}.
\item
  {[}BH{]} K. J. Böröczky and M. Henk, \emph{Cone-volume measure and
  stability}, Advances in Mathematics 306 (2017), 24--50, DOI
  10.1016/j.aim.2016.10.005.
  \href{https://users.renyi.hu/~carlos/cone-volume-stability.pdf}{Author-hosted
  text}.
\item
  {[}S{]} C. Saroglou, \emph{Volumes of projection bodies of some
  classes of convex bodies}, Mathematika 57 (2011), 329--353, DOI
  \href{https://doi.org/10.1112/S0025579311001860}{10.1112/S0025579311001860}.
\item
  {[}Sch{]} R. Schneider, \emph{Convex Bodies: The Brunn--Minkowski
  Theory}, second expanded edition, Cambridge University Press, 2014.
  Standard support-function, surface-area, projection and mixed-volume
  inputs.
\end{itemize}

Versions 2--4 preserve and credit the product and simplex computations
inherited from OpenAI family 088, pinned at commit
\texttt{adc7f1241b42e322a6451854ab7e4b4c146bf78a}. This supplement adds
an explicit quantitative proof to that continuing record; it does not
replace historical versions. The manuscript is AI-assisted research and
makes no first-discovery assertion.

\end{document}
